\chapter*{Abstract}
\addcontentsline{toc}{chapter}{Abstract}

Over ninety AI models have been built to predict sepsis in real time, yet they
fail in practice. Wang et al.\ (2025) showed that while AUROC barely changes
between internal and external testing (0.811 to 0.783), the clinical Utility
Score collapses from $+0.381$ to $-0.164$, meaning these models cause more
harm than benefit when deployed. This thesis investigates why: is the problem
the model, the sepsis label, or the researcher's methodological choices?

The study uses MIMIC-IV v3.1 (65,078 adult ICU stays) with external validation
on eICU-CRD v2.0 (63 sepsis cases across 7.9 million person-hours). Three
Sepsis-3 label variants (Narrow, Seymour-Standard, Liberal) are tested with a
dynamic discrete-time competing-risks hazard model, gradient-boosted trees (GBT),
a static Cox model, and rule-based scores (NEWS2, qSOFA, SIRS). Because all
three variants share one culture-plus-antibiotic anchor, the Sepsis-3
conjunction is additionally split into its two limbs and each is labelled
separately, including a deterioration label that uses no treatment timestamp at
all.

The results show that the sepsis label definition changes AUROC by 0.153
(95\% CI 0.136--0.177), a spread statistically indistinguishable from the one
produced by changing the model (0.142, 95\% CI 0.130--0.161): both analyst
decisions are large and neither dominates. Ten of the twenty shared clinical
covariates reverse direction across labels, making published risk factors
unreliable without knowing which label was used. Random data splits inflate
AUROC by 0.123 (95\% CI 0.100--0.147) compared to time-based splits. Most
importantly, the pre-treatment evaluation window contains \textbf{zero sepsis
events} in 179,003 person-hours: no patient develops sepsis before doctors
begin treatment. This means current models predict when clinicians act, not
when patients deteriorate. The Utility Score is \textbf{negative for every
model and every label}, with over 90 false alerts for each true one.

Splitting the definition shows how much of the label effect is definitional and
how much is real. The onset \emph{time} under every Sepsis-3 variant is fixed
entirely by the treatment anchor, the organ-dysfunction criterion selecting
which stays are labelled without moving any onset by an hour. Yet a
treatment-independent deterioration label agrees with Sepsis-3 at
$\kappa = 0.057$, close to chance, so the label effect is not an artefact of
re-reading one definition. Sepsis-3 also proves \emph{harder} to predict than
either limb it is built from --- 0.607 AUROC against 0.756 for the treatment
anchor alone and 0.842 for deterioration alone --- which places the difficulty
in the conjunction rather than in the physiology or in clinician behaviour.
Critically, the same deterioration label records \textbf{1,362 onsets in the
pre-treatment window where Sepsis-3 records none}, at 0.840 AUROC: patients do
deteriorate before clinicians act, and the empty window is a limitation of the
definition rather than a fact about the patients.

All estimates carry stay-level cluster-bootstrap intervals; the six
pre-registered hypotheses are corrected by Holm--Bonferroni and all subgroup
analyses by Benjamini--Hochberg. Under correction, subgroup discrimination
differences by race are too imprecisely estimated to support any claim, while
differential \emph{label} sensitivity does survive: patients whose race is
undocumented, and Spanish-speaking patients, change sepsis status across label
definitions significantly more often than the reference groups.

The collapse in clinical utility is not a model problem; it is a label
problem. Better models cannot fix a label that depends on the actions it is
supposed to predict ahead of.

\noindent\textbf{Keywords:} sepsis; MIMIC-IV; discrete-time survival analysis;
competing risks; label definition; variance decomposition; clinical utility;
temporal dataset shift; treatment-anchored labels; TRIPOD+AI.
